\section{Introduction}

ChatGPT was launched by OpenAI november 3oth of 2022. It is powered by GPT-3. A \glsxtrfull{LLM} trained wiht massive data across WebText2. This includes websites, books, articles and toher types of text based information. It is cappable of analyzing stadistical relationships between words and phrases. This way, the model is extraordinarily efficient in tasks such as generating, summarizing and translating text in a coherent way. \citep{RAY2023121}

Even though many of the capabilities of ChatGPT were already maneagable with GPT-2 since 2019, the main different aspect of GPT-3 which I will be expLoiting is its universal accessibility. In november 2022 the launch of ChatGPT stormed into society's psyche via public use with no payment requiered, an easy interface simple enough to only input a question and an easygoing nature capable of running on most devices. In january 2023, 2 months after its launch, the agent reached around 100 million users. \footnote{\cite{paris2023chatgpt}} 

The fast technological shift due to \glsxtrshort{LLM}'s has sparked a heated debate as enterprises and governments adapt to the new paradigm. This debate is central to the thesis I present as one of the main points of disagreement is the effect that \glsxtrshort{LLM}'s have on the labor market.

In economic theory the effects are ambiguous. On one hand, \glsxtrshort{LLM}'s can be mere complements by enhancing worker's performance. ***** on the other hand, it completely subsitutes out workers in certain fields. \citep{hui2024short} Thus it is of great importance to the mexican market to analyze the function an \glsxtrshort{LLM}, such as ChatGPT, has on tasks in order to understand potential wage differentials and job transitions. 

Previous literature has found that adoption of other similar tecnhologies such as computers can worsen wage inequality and affect negatively employment levels. \citep{krueger} \citep{katz} \citep{autor2003} Moreover, the mexican context is a particular case in which its persistent informality can weaken or strengthen these effects. 

With the launch of ChatGPT as an exogeonous event I use the rotative panel \glsxtrfull{ENOE} and estimate the emplyment and wage effects in the mexican labor market. My main specification is a \glsxtrfull{DID} and it is inspired by \cite{hui2024short}'s model. I classify exposed tasks and construct an occupation\footnote{I construct an individual index for all occupations in \glsxtrshort{ENOE} according to \glsxtrfull{SINCO}} specific index relying on methodology by \cite{martinez2024risk} and \cite{eloundou2023gptsgptsearlylook}.

I provide the first estimates for labour and wage impacts due to the launch of ChatGPT in Mexico. These reults reveal XXXX.


%% Añadir analisis de compsicion industrial 
